\patchcmd{\chapter}{\thispagestyle{plain}}{\thispagestyle{fancy}}{}

\chapter{INTRODUCTION}
\patchcmd{\chapter}{\thispagestyle{plain}}{\thispagestyle{fancy}}{}

\hspace{0.5 cm}
The rise in software system's development and the rising complexity in application development processes has made the need for good code documentation necessary.  Software systems tend to become more advanced in terms of features, connections and development teams. Code documentation helps make code readable, encourages teamwork between programmers and guarantees software system maintainability in the future. Documentation acts as a link between programmers by offering an explanation of system features, design choices and coding processes. In large projects involving many programmers working on the project for a long time, documentation is vital to ensuring consistency and minimizing misunderstandings. Nonetheless, keeping documentation up to date and accurate is difficult because there are many modifications to code and it is time-consuming to write documentation manually.
\par Often, in practice, in many software development projects, it is observed that documentation tends to fall behind and there are chances that the documentation might not be complete for the software developed. In such situations, the developers do not have enough time to document the changes made to the codebase, owing to their focus being diverted to implementing features and fixing bugs. It also results in an inconsistency in the actual state of the codebase and what is described in the documentations related to the code. This makes it challenging for developers to understand the behavior of the system and hence spend extra time debugging and analyzing the actual code.
\newpage
\par For addressing these difficulties, a set of latest development of Artificial Intelligence (AI) technologies can be used to automate documentation creation. The analytical technique of Abstract Syntax Tree (AST) is one of the most useful ways of analyzing source codes by building an abstract syntactic tree. Such components include functions, classes, variables and control statements, thus helping to have a clear understanding of the structure of source codes. Compared to traditional text approaches, analysis through AST allows to get more precise results and insights on behavior of the source codes analyzed. 

\par Apart from the AST-based analysis, another way of gaining deeper insight into code by analyzing relationships among its different elements is by means of graphical representation. In this method, nodes are represented using objects such as functions, classes and modules, while edges are represented using relationships between them, such as function call dependencies and inheritance. Graphs can be very useful to identify relationships and patterns in the form of dependencies, control flow and the most important components in the system. So, together with the AST analysis method, graphs give a broader perspective on the code base. 

\par In addition, Large Language Models (LLMs) have brought another evolution in the documentation process, where the complex and complicated structure of code can be translated into understandable and readable text. The detailed grasp of the semantics of the code by LLMs is useful in making fine descriptions, summaries and documentation of the code itself. Using structured input based on the analysis of an abstract syntax tree and graphs, LLMs can provide precise, consistent and contextually relevant documentation, README files and architectural documentation of modules and functions.

\par The goal of this project is to create an intelligent system for automatic generation of documentation based on source code. This will be achieved through the use of AST-based parsing and modeling techniques, as well as natural language models, to reduce human interference, enhance software readability and maintainability. This intelligent system will aid developers in creating accurate and updated documentation without having to modify the source code.

\par The increasing demand for intelligent and scalable documentation highlights the limitations of conventional methods in handling complex software systems. Manual documentation processes are often slow, inconsistent and unable to keep up with rapid development cycles and frequent software updates. As a result, maintaining accurate and up-to-date documentation becomes a significant challenge.

\par In contrast, integrating structural analysis techniques with state-of-the-art language models offers a more effective solution by enabling automated and context-aware documentation generation. Therefore, reviewing the existing literature is essential to identify current approaches, their limitations and potential areas for improvement, which are further explored in the following chapter.